\chapter{Polishing the Causal Stream: A Pluralistic Reading of the \emph{Platform Sutra} and a Parfitian critique of the traditional interpretation}
\chaptermark{A Pluralistic Reading of the \emph{Platform Sutra}}
\chapterauthor{Ng Kang Qi,
\textit{University of Cambridge}}

\renewcommand*{\thesection}{\arabic{section}.}
\renewcommand*{\thesubsection}{\arabic{section}.\arabic{subsection}.}

\begin{quote}
    \emph{The} \emph{Platform Sutra of the Sixth Patriarch} is a pivotal text in Chinese Buddhist philosophy, traditionally interpreted as establishing the orthodoxy of the Southern ``Sudden" school of Chan Buddhism while relegating the Northern ``Gradual" school to a provisional teaching. To understand this schism, one must understand the mechanics of Buddhist soteriology: enlightenment fundamentally requires liberation from the ``psychological stream"---the continuous, causally linked flow of mental states, habits, and karmic dispositions that perpetuates the illusion of a permanent self. This paper argues that the traditional interpretation fails to properly address this causal reality by making a fatal logical leap from necessity to sufficiency. By claiming that a sudden, propositional insight into the emptiness of the self is entirely sufficient for enlightenment, the orthodoxy invites antinomianism and dogmatically dismisses the Northern method of cultivation. Drawing upon Derek Parfit's reductionist metaphysics and Gilbert Ryle's epistemology, I argue that the Northern school's gradual approach does not mistakenly reify a Cartesian ego but is rather strictly necessary to convert the propositional truth of emptiness (knowing-that) into the dispositional purity of enlightened living (knowing-how). Ultimately, this paper rehabilitates the \emph{Platform Sutra} as espousing a methodological pluralism: true liberation requires both the sudden, conceptual destruction of the ego and the continuous, gradual cultivation of one's dispositions.
\end{quote}

\vspace{\credgap}

\section{Introduction}

In Chinese Buddhist philosophy, the \emph{Platform Sutra of the Sixth Patriarch} is regarded as the definitive manifesto of the ``Southern
School'' of Chan Buddhism. The ``Southern School'' represents an alleged
doctrinal divergence from the ``Northern School'', with both branches
being rooted in a historical Mahayana tradition descending from the
shared ``East Mountain'' teachings of the Fifth Patriarch, Hongren.\footnote{John R. McRae, \emph{The Northern School and the
  Formation of Early Ch'an Buddhism} (Honolulu: University of Hawaii
  Press, 1986), 30--34. (Note: McRae thoroughly details Hongren's ``East
  Mountain Teaching'' as the shared progenitor of both schools).} The doctrines coincide in their understanding of enlightenment as a
liberation from \emph{Samsara}---the cycle of life and death---through a detachment from worldly desire. However, both schools diverge in their
interpretations of what this detachment consists of, and how it can be achieved.

According to the traditional interpretation as established by Southern
polemicists, the most salient soteriological differences between these
schools can be represented in both their views of the nature of
enlightenment and the method through which enlightenment is
attained\footnote{The traditional interpretation can be traced primarily
  to the rhetoric of Heze Shenhui (684--758), a disciple of Huineng. In
  his \emph{Treatise Establishing the True and False according to the
  Southern School of Bodhidharma}, Shenhui explicitly attacked Shenxiu's
  lineage for advocating ``gradual'' cultivation. He condemned their
  practices---such as ``freezing the mind to enter concentration'' and
``fixing the mind to view purity''---as dualistic errors that failed to
  apprehend the sudden, unconditioned nature of enlightenment. See John
  R. McRae, \emph{Seeing through Zen} (Berkeley: University of
  California Press, 2003), 33--35.}:

\vspace{-1.3em}
\begin{table}[h!]
    \caption{The Traditional Interpretation's Characterisation of the
Soteriological Divide.}
\vspace{0.5em}
    \label{tab:table1}
    \begin{tabular}{|p{0.23\linewidth}|p{0.33\linewidth}|p{0.33\linewidth}|} % <-- Alignments: 1st column left, 2nd middle and 3rd right, with vertical lines in between
      \hline
       & \textbf{Nature of Enlightenment} & \textbf{Method of Attainment}\\
      \hline
      Northern School & A dispositional and recurrent state of purified moral and cognitive habits; the continuous maintenance of an unattached mind, free from worldly striving, which ultimately secures liberation from samsara and the cessation of suffering. & Gradual character cultivation, sustained meditation, and continuous moral vigilance to prevent moral defilement. \\
      \hline
      Southern School & An occurrent, transformative realisation of a propositional truth: the ontological emptiness of the self and the inherent purity of the mind, which directly results in liberation from samsara and the cessation of suffering. & Direct, unmediated intuitive awakening; a sudden paradigm shift that bypasses conceptual study. \\
      \hline
    \end{tabular}
\end{table}

\noindent According to the traditional interpretation, the doctrinal divergence
illustrated above is the main rhetorical angle of the \emph{Platform
Sutra} with the text dramatising this schism through a pivotal narrative
event: a succession contest wherein the learned head monk Shenxiu and
the illiterate layman Huineng are framed as rhetorical opposites. While
Shenxiu's submitted verse advocates for gradual character cultivation,
Huineng's verse characterises enlightenment as a sudden flash of insight
into a reductionist ontology of the self. Within the narrative of the
\emph{Platform Sutra}, the outcome is decisive: the Fifth Patriarch
acknowledges Huineng's verse as depicting the ultimate truth and crowns
him the Sixth Patriarch.

Building upon this narrative outcome, the traditional interpretation
asserts three conclusions. First, regarding the \emph{nature} of
enlightenment, it concludes that an occurrent realisation of Huineng's
metaphysics is entirely sufficient for liberation. Second, it dismisses
Shenxiu's views entirely, claiming his framework of gradual cultivation
is an elementary metaphysical error that reifies a non-reductionist
self. Third, regarding the \emph{method} of enlightenment, it asserts
that the only means through which liberation can be attained is by
sudden insight.

The objective of this paper is to challenge these interpretative
conclusions by arguing that the traditional interpretation overplays the
extent to which the narrative victory of Huineng should represent a
corresponding dismissal of Shenxiu. While the \emph{Platform Sutra}
places great emphasis on the importance of ``Sudden Enlightenment'' and
a reductionist ontology of the self, I argue that Shenxiu's views on the
nature of enlightenment as a disposition attained through gradual
cultivation are compatible with such an ontology. Ultimately, I claim
that Huineng's soteriology must necessarily contain some aspects of
Shenxiu's teachings to be tenable. To demonstrate this, we must
re-evaluate the implications of the traditional interpretation and show
that the three main soteriological conclusions it draws are unjustified.

The paper proceeds as follows. In Section 2, I argue against the
traditional interpretation's claim that propositional knowledge is
sufficient for enlightenment by showing that this view leads to
antinomianism, and thus cannot be an accurate reflection of the
\emph{Platform Sutra}. Consequently, the Northern view of dispositional
purity is necessary to a tenable Chan soteriology. In Section 3, drawing
on Derek Parfit, I demonstrate that the traditional interpretation
misinterprets Shenxiu's ontology, and that the pedagogy of cultivating
dispositions need not reify a non-reductionist self. In Section 4,
drawing from Gilbert Ryle's thesis regarding the independence of
propositional knowledge and dispositional skill, I argue against the
traditional interpretation's view of the method of attaining
enlightenment; because sudden insight is insufficient to build
dispositional purity, the Shenxiu's teachings of gradual cultivation are
strictly necessary to a tenable pedagogy of working towards
enlightenment. Finally, in Section 5, I elaborate on why Huineng's
propositional insight remains a necessary prerequisite for this gradual
cultivation. This clarifies the narrative logic of the \emph{Platform
Sutra}, justifying the Fifth Patriarch's prioritisation of Huineng while
exposing the root of the traditional interpretation's error: conflating
necessity with sufficiency. Ultimately, I propose that the
\emph{Platform Sutra} should be viewed as espousing a pluralistic theory
of enlightenment: one that defines the nature of enlightenment as both
propositional insight and dispositional purity, and the method of
attainment as both sudden awakening and gradual cultivation. This
synthesis, corroborated by Guifeng Zongmi, seamlessly integrates the
doctrines that the traditional interpretation falsely forces into a
mutually exclusive ``Northern" and ``Southern" dichotomy.

\section{The key verses of the Platform Sutra}

The philosophical core of the \emph{Platform Sutra} and the source of
the tension addressed in this paper is found in the first section of the
Sutra, where two verses are presented at the behest of the Fifth
Patriarch---the great master of the Baolin monastery---as a means of
finding a successor. According to the Sutra, the verse contest would
allow the Fifth Patriarch to determine which disciple had apprehended
the deeper truth of enlightenment and could therefore lead the monastery
towards an appropriate soteriology.\footnote{Philip B. Yampolsky,
  trans., \emph{The Platform Sutra of the Sixth Patriarch} (New York:
  Columbia University Press, 1967), 128--129.}

The first verse, presented by the learned head monk Shenxiu, is a
representation of the ``Northern'' school of Chan:

\begin{CJK*}{UTF8}{gbsn}
\begin{quote}
The body is a bodhi tree (\emph{shēn shì pú tí shù}, 身是菩提树),

The mind is like a bright mirror stand (\emph{xīn rú míng jìng tái},
心如明镜台),

At all times we should diligently polish it (\emph{shí shí qín fú shì},
时时勤拂拭),

and prevent dust from collecting (\emph{wù shǐ rě chén āi},
勿使惹尘埃).\footnote{Yampolsky, \emph{The Platform Sutra}, 130}
\end{quote}
\end{CJK*}

\noindent Within the narrative, Shenxiu's verse is treated as an incomplete
apprehension of the true nature of enlightenment, and he is described by
the Patriarch as ``unable to see the self-nature,'' suggesting his
framework is provisional rather than ultimate. After Shenxiu, the
illiterate layman Huineng's verse is revealed:

\begin{CJK*}{UTF8}{gbsn}
\begin{quote}
Bodhi originally has no tree (\emph{pú tí běn wú shù}, 菩提本无树),

The bright mirror also has no stand (\emph{míng jìng yì fēi tái},
明镜亦非台),

Fundamentally, there is not a single thing (\emph{běn lái wú yī wù},
本来无一物),

Where could dust arise? (\emph{hé chù rě chén āi},
何处惹尘埃)\footnote{Yampolsky, \emph{The Platform Sutra}, 132.}
\end{quote}
\end{CJK*}

\noindent In contrast to Shenxiu, Huineng is celebrated by the Fifth Patriarch in
private as having truly perceived the fundamental, unconditioned nature
of the mind, thereby earning the secret transmission of the robe and
Dharma, and the title of the Sixth Patriarch.

\subsection{The traditional interpretation}

According to the traditional interpretation initially published by
Southern polemicist Shenhui and later accepted by standard Chan
historiography, the reason for Huineng's rhetorical victory lies in how
these verses map onto a fundamental ontological dispute. In Western
analytic terms, this debate directly parallels Derek Parfit's
distinction between non-reductionist and reductionist metaphysics of the
self.

According to this view, Shenxiu's verse espouses a non-reductionist
metaphysics of the self. Because the verse makes a distinction between
the mind (``Mirror Stand") and its transient attachments or impurities
(``Dust"), it is claimed that Shenxiu implies that the mind is a
container that exists independently of its contents. In modern
metaphysical terms, Shenxiu is accused of positing a ``Further Fact" of
identity: a substantive, unchanging self that exists over and above the
stream of experience.

However, this substantive self is incompatible with the fundamental
Buddhist doctrine of \emph{Anicca} (impermanence)---because
\emph{Anicca} posits that all phenomena are in a state of constant,
interdependent flux, the existence of a permanent, unchanging, and
independent essence or soul must be rejected.\footnote{Mark Siderits,
  \emph{Buddhism as Philosophy: An Introduction} (Indianapolis: Hackett
  Publishing, 2007), 32--34.} The traditional orthodoxy thus uses this
rhetoric as a means of discrediting Shenxiu's view. According to this
interpretation, because Shenxiu commits a fundamental metaphysical
error, his practical teachings must be dismissed as a category mistake.
If there is no permanent, underlying entity to accumulate defilement,
the act of continuous purification is logically incoherent, since one
cannot polish a non-existent `mirror'. Therefore, the gradual method is
tolerated within the Southern narrative only as a remedial crutch
designed solely to prevent practitioners with `dull faculties' from
falling into grave moral error.

The traditional narrative then frames Huineng's verse as the welcome
reductionist reply: by asserting ``The mirror has no stand," Huineng
denies the existence of some `further fact' about the self, moving away
from the problematic non-reductionism that Shenxiu apparently posits.
Crucially, the traditional orthodoxy goes even further: it concludes
that because Huineng is right about the ontology of the self, the nature
of enlightenment is constituted merely by a propositional understanding
of this ontology.

This conclusion is problematic for two primary reasons. First, it uses
Huineng's rhetorical victory to justify a dangerous view of
enlightenment that is likely to lead to antinomianism. It assumes that
because Huineng holds the correct reductionist view of the self, an
occurrent understanding of this metaphysical truth is inherently
sufficient for enlightenment. Second, this narrative is extremely
uncharitable to Shenxiu, prematurely reducing his gradualist pedagogy to
an elementary metaphysical mistake.

\subsection{Semantic ambiguity and Huineng's ontology}

To demonstrate why treating this propositional knowledge as sufficient
is fatal to Chan soteriology and precisely how it invites the threat of
antinomianism, we must first clarify the exact reality Huineng is
describing. This requires a close analysis of his main metaphysical
claim: ``fundamentally, there is not a single thing'' (\emph{běn lái wú
yī wù}, \begin{CJK*}{UTF8}{gbsn}本来无一物 \end{CJK*}). The main challenge in understanding this assertion
lies in the key term ``a single thing'', which is semantically
overloaded. Depending on how we interpret this phrase, the verse
generates three entirely distinct and mutually exclusive ontological
claims:

\vspace{0.5em}
\noindent Reading 1: The Nihilistic Reading (``Single thing"= Anything)

\vspace{0.5em}
\noindent If ``a single thing" is read as `anything', then the statement denies the
existence of any objects (variables) whatsoever within the domain.

\vspace{0.5em}
\noindent $\neg\exists x(x=x)$

\vspace{1.5em}
\noindent Reading 2: The Anti-Monistic Reading (``Single thing"= One thing)

\vspace{0.5em}
\noindent If the word ``single" is read instead as a numerical quantifier
designating one, then the statement denies that there is exactly one
object within the domain (meaning there is either zero or more than one
object).

\vspace{0.5em}
\noindent $\neg(\exists x(x=x) \land \forall x \forall y(x=y))$

\vspace{1.5em}
\noindent Reading 3: The Anti-Pluralistic Reading (``Single thing"= Distinct
thing)

\vspace{0.5em}
\noindent If ``a single thing" is read as a substantive, independent entity
possessing inherent existence, then the statement denies the existence
of boundaries, or of distinctness between objects. It asserts that there
are no distinct things in the universe.

\vspace{0.5em}
\noindent $\neg \exists x \exists y(x \neq y)$

\vspace{1.5em}
\noindent Which of these readings is the most plausible? We can immediately reject
Reading 1 (Nihilism), as it is directly incompatible with the rest of
Huineng's soteriology. An empty universe precludes the possibility of
the experiencing mind and sudden awakening, which are the foundations of
Huineng's brand of Chan Buddhism. Furthermore, elsewhere in the
\emph{Platform Sutra}, Huineng explicitly states:

\begin{quote}
    ``Emptiness includes the sun, moon, stars, and planets, the great earth,
mountains and rivers, all trees and grasses, bad men and good men, bad
things and good things, heaven and hell; they are all in the midst of
emptiness."\footnote{Yampolsky, \emph{The Platform Sutra}, 146.}
\end{quote}


\noindent This passage serves a crucial philosophical purpose: it definitively
populates the ontology. Because the nature of emptiness encompasses the
phenomenal world, the domain of discourse cannot be empty, and the
nihilistic reading must be rejected. Furthermore, the rhetorical
structure of this passage eliminates Reading 2 (Anti-Monism). Huineng
deliberately lists the vast pluralities and dualities of the phenomenal
world, including the good and bad, heaven and hell, only to immediately
collapse them into the singular concept of ``emptiness." From this we can
surmise that he is not advocating for pluralism but trying to dismantle
it.

Therefore, we are left with Reading 3. When Huineng asserts ``there is
not a single thing," he is using the formulation $\neg \exists x \exists y(x \neq y)$ to make a
mereological claim that denies the existence of boundaries and
distinctness. In standard logic, the statement ``no two things are
distinct" is logically equivalent to the universal affirmation of
identity: $\forall x \forall y(x=y)$ (everything is identical to everything else). When
we combine this anti-pluralistic premise ($\forall x \forall y(x=y)$) with our rejection
of the nihilistic reading, the conclusion is that Huineng is asserting
that there exists \emph{exactly one} thing $\exists x(x=x \land \forall y(y=x))$.

This conclusion that reality is an undifferentiated whole is grounded in
the broader textual tradition of Chan Buddhism. In the \emph{Platform
Sutra} and wider Mahayana literature, this singular underlying reality
is explicitly identified as the inherently pure Buddha-nature (\emph{fó
xìng}, \begin{CJK*}{UTF8}{gbsn}佛性\end{CJK*}). The denial of pluralistic boundaries is the ontological
equivalent of the Chan doctrine of non-plurality, which posits that
there is no substantive metaphysical division between ordinary sentient
beings and the enlightened Buddha.

\subsection{Monism and Antinomianism}

This brings us to the core tension. While Huineng's metaphysical monism
aligns with the broader Mahayana doctrine of Buddha-nature, the
traditional interpretation weaponises this ontology to make an
unjustified conclusion regarding the nature of enlightenment. The
orthodoxy assumes that because reality is an inherently pure,
unpartitioned whole, the mere propositional realisation of this
metaphysical truth is sufficient to \emph{constitute} enlightenment.

However, this cannot be the case. Treating propositional knowledge as
the sole requirement for enlightenment invariably leads to
antinomianism: the dangerous theological doctrine that moral laws,
ethical habituation, and behavioural precepts are no longer binding. If
enlightenment is reduced to a sudden flash of conceptual insight into
the underlying purity of reality, it implies there is nothing left to
do, and that the everyday Buddhist practice of continuous moral
cultivation is rendered useless. Under this view, one could grasp the
monistic ontology while continuing to engage in harmful behaviour,
justifying their actions under the guise that ``all things are already
the Buddha-nature."

This antinomian outcome is unacceptable for two reasons. First, it is
practically disastrous for the Chan tradition, which is premised on the
idea that Buddhist practitioners must aim to alleviate real-world
suffering. Second, it would be unable to account for other sections of
the \emph{Platform Sutra}.

For example, in Chapter 6, Huineng explicitly confers the ``Formless
Repentance" (\emph{wú xiàng chàn huǐ}, \begin{CJK*}{UTF8}{gbsn}无相忏悔\end{CJK*}) upon his disciples,
commanding them to actively repent for past and future sins arising from
ignorance and deceit. Crucially, he does not merely ask them to realise
the abstract emptiness of these sins; he explicitly demands that they
``awaken to {[}their{]} own errors and \textbf{reform} {[}their{]} past
evil."\footnote{Yampolsky, \emph{The Platform Sutra}, 143--144.}

This pedagogical mandate refutes the traditional interpretation's claim
that a sudden flash of propositional understanding is sufficient to
constitute enlightenment. The act of ``{[}reforming{]} past evil"
inherently requires a functional, temporal separation between past
karmic defilement and present moral action. It demands a continuous
ongoing change in psychological habituation and behaviour, beyond a
change in metaphysical understanding. If a practitioner only possesses
the propositional knowledge that everything is one with the
Buddha-nature, yet remains acting in a manner that brings more suffering
to the world, they cannot be enlightened according to the doctrines of
Chan. Therefore, the nature of enlightenment must necessarily include a
dispositional purity alongside propositional insight.

If this dispositional purity is practically necessary, as Huineng's own
repentance rituals demonstrate, then the dismissal of Shenxiu's verse is
premature. Shenxiu's instruction to ``diligently polish" the mirror
represents the exact conceptual framework required to execute the
behavioural changes that the traditional interpretation abandons. To
save Chan practice from the complacency of the absolute, a more
charitable interpretation of Shenxiu's gradualism is in order. We must
reinterpret his verse not as an incorrect metaphysical claim, but as a
practical and behavioural framework for maintaining dispositional
purity.

This rehabilitation, however, presents a significant philosophical
hurdle: how can this be done without risking a problematic,
non-reductionist ontology of the self that Shenxiu seems to imply? To
resolve this tension, we require a philosophical framework capable of
distinguishing between the metaphysical subject (which Huineng rightly
targets) and the psychological continuum (which the traditional reading
inadvertently discards).

\section{The Parfitian Solution}

Derek Parfit's reductionist metaphysics in \emph{Reasons and Persons}
provides precisely the analytical vocabulary needed to resolve this
tension. Parfit allows us to strictly separate the question of identity
(the metaphysical subject) from the question of survival (the
psychological continuum), revealing that the ``Sudden" and ``Gradual"
approaches are not mutually exclusive, but rather address two entirely
distinct ontological categories.\footnote{Derek Parfit, \emph{Reasons
  and Persons} (Oxford: Oxford University Press, 1984), 215.}

To map Parfit's framework onto the Chan debate, we must first understand
his distinction between two rival views of persons:

\vspace{0.5em}
\noindent Non-reductionism (The Cartesian View): this view holds that personal
identity is ensured by a ``Further Fact", or a separately existing
entity, such as a Cartesian Ego or a Soul, that remains distinct from
the brain and body.

\vspace{0.5em}
\noindent Reductionism (The Bundle View): Parfit, aligning with Hume and orthodox
Buddhism, holds that a person is not a separate underlying entity, but a
composite bundle of interrelated mental and physical events. For the
Reductionist, personal identity is constituted solely by Relation R:
psychological connectedness and continuity (memory, character, and
causal links).\footnote{Parfit, \emph{Reasons and Persons}, 216--217.
  See also David Hume, \emph{A Treatise of Human Nature}, 1.4.6.}

\vspace{0.5em}
\noindent Applying these categories provides the necessary lens to interpret
Shenxiu without uncharitably reducing his view to non-reductionism.
Rather than thinking of Shenxiu's verse as reifying a ``Further Fact''
about the self, we can interpret his verse as a necessary pedagogy
regarding Relation R.

\subsection{Identity vs. Survival: The Independence of Relation R}

To understand why this pedagogy regarding Relation R remains practically
necessary, we must utilise Parfit's distinction between identity and
survival. Parfit argues that what matters in personal existence is not
the persistence of a ``Further Fact" (identity), but the preservation of
relation R (survival).\footnote{Parfit, \emph{Reasons and Persons}, 255.
  Parfit argues that ``Identity is not what matters in survival."} The
traditional interpretation correctly identifies Huineng's verse as a
metaphysical attack on the ``Further Fact." Huineng's assertion that
there is ``not a single thing" simultaneously asserts a Buddhist monism
and attacks the idea of the Cartesian Ego, aligning with the Buddhist
doctrine of \emph{anatman} (no-self). However, the dismissal of Shenxiu
leads directly to antinomianism because it overlooks a crucial reality:
the negation of the ``Further Fact" does not negate Relation R. Even in a
strictly reductionist universe where there is no permanent self to `own'
experiences, the causal stream of psychological continuity continues to
flow. Mental state M1 (e.g., a moment of anger) remains causally
connected to mental state M2 (e.g., a disposition toward violence) via
memory and psychophysical causality. This causal chain exists
independently of the ontology of the ``Further Fact." Ignoring the
necessity of regulating Relation R is what leads the traditional
interpretation towards antinomianism.

Therefore, when Shenxiu commands practitioners to ``diligently polish"
the mirror, we should not read this as the reification of a
non-reductionist view of the self. Rather, it is the necessary
pedagogical mechanism for managing Relation R, where Shenxiu is
addressing the causal reality of the psychological continuum. Just
because the ``polisher" and the ``mirror" lack an independent, substantial
existence does not mean the causal links between immoral thoughts and
harmful actions cease to operate. Gradual cultivation is the essential
practice of purifying the causal stream, and it is the only way to make
sense of Huineng's command for ``formless repentance" and avoid
antinomianism.

One might object that the syntactical structure of Shenxiu's verse
contradicts this interpretation and commits him to a non-reductionist
ontology. By utilising nouns like ``mirror stand" and ``dust," in a
subject-copula-predicate form (``The body \textbf{is} a bodhi tree''),
he appears to be hypostatising the mind, treating it as a substantive
container. However, a closer reading of the original Chinese text
reveals this to be a functional analogy rather than a strict identity
statement. While his first line uses the copula to assert ``The body
\emph{is} {[}\emph{shì}, \begin{CJK*}{UTF8}{gbsn}是\end{CJK*}{]} a bodhi tree," his second line
deliberately shifts to a simile: ``The mind is \emph{like} {[}\emph{rú},
\begin{CJK*}{UTF8}{gbsn}如\end{CJK*}{]} a bright mirror stand." From this, we can grant that Shenxiu is
not making a non-reductionist claim about the self but is instead
utilising an analogy to highlight the mind's capacity to reflect reality
clearly or be obscured by passing phenomena.

Furthermore, the syntactical phrasing is likely a byproduct of
pedagogical necessity. Language is inherently dualistic; the very syntax
of a behavioural imperative requires a grammatical subject and object.
To describe the stewardship of the mental stream without using the
convenient shorthand of a ``mirror" and ``dust" would require a
cumbersome, technical circumlocution that would be pedagogically inert
and perhaps even more opaque to the average practitioner.

Thus, Shenxiu adopts a pragmatic ``intentional stance." Just as a modern
biologist might say ``the gene wants to replicate" without literally
ascribing conscious desire to a molecule, Shenxiu commands his monks to
``polish the mirror" without literally ascribing a permanent Cartesian
soul to the practitioner. He is utilising conventional language to point
toward necessary practical action. Once we strip away this grammatical
hypostatisation, we can accurately reconstruct the underlying logical
form of Shenxiu's verse. When he commands his practitioners to
``diligently polish" and ``prevent dust from collecting," he is
fundamentally asserting that there is a causal sequence of mental states
that requires continuous regulation.

In Parfitian terms, we can map his metaphor directly onto the mechanics
of Relation R. The ``dust" is not a physical property attached to an
underlying soul, but rather the causally salient factors within the
psychological continuum, such as fleeting thoughts of greed, anger, or
ignorance. If the person at time T1 allows this ``dust" to settle by
reinforcing a habit of ignorance, the person at time T2 will inherit
those dispositions via Relation R. ``Polishing," therefore, represents
the precept of regulating a particular causal stream so that future
psychological states develop the correct moral dispositions. The
traditional interpretation, by fixating exclusively on Huineng's
negation of the ``Mirror," throws out the baby (Relation R) with the
bathwater (the ``Further Fact"). It leaves the practitioner with a
metaphysical theory, but without a practical method to address the
continuity of their conditioned habits. By utilising Parfit, we
successfully vindicate Shenxiu's verse as necessary framework for
understanding the dispositional nature of enlightenment.

\section{The epistemology of enlightenment with lessons from Ryle}

Having established that the nature of enlightenment inherently requires
dispositional purity (the purification of Relation R), we must now turn
to the method of attainment. A defender of the traditional
interpretation might concede that the nature of enlightenment
necessarily requires a dispositional aspect, but resist the idea that
the method of attaining enlightenment requires gradual cultivation by
contending that a sufficiently profound, sudden propositional insight
would automatically generate these dispositions.

However, drawing upon the epistemology of Ryle reveals why this argument
must fail. Ryle would call such a view the ``intellectualist legend": the
false assumption that intelligent, skilful action is merely a derivative
of theoretical knowledge, or that to ``know" a truth is to automatically
be able to act upon it. In \emph{The Concept of Mind}, Ryle argues
against the view that performance is merely the execution of a
recognised rule. He distinguishes between two cognitive categories:

\vspace{0.5em}
\noindent Knowing-That (Propositional Knowledge): the possession of information or
theoretical truth (e.g., knowing \emph{that} a bicycle functions via
gyroscopic force).

\vspace{0.5em}
\noindent Knowing-How (Dispositional Skill): the capacity to perform a task or
navigate a domain (e.g., knowing \emph{how} to ride a
bicycle).\footnote{Gilbert Ryle, \emph{The Concept of Mind} (London:
  Hutchinson, 1949), 27--32.}

\vspace{0.5em}
\noindent Crucially, Ryle argues that these two categories are fundamentally
distinct and irreducible: one can possess propositional knowledge in a
subject matter but lack the corresponding dispositional skills, and vice
versa. Applying this to the Chan debate reveals a crucial point. If one
accepts that the nature of enlightenment requires the conditioning of
one's causal stream (which the preceding sections have established is
necessary to avoid antinomianism), then mere propositional knowledge of
emptiness will not automatically manifest this dispositional state.
Huineng and Shenxiu thus address two distinct epistemological categories
of the same soteriological goal. Huineng's verse concerns the
propositional knowledge of the nature of the self in relation to the
universe, while Shenxiu is espousing the pedagogy of the dispositional
skill of living according to this ontological reality. The error of the
methodological monism promoted by the traditional interpretation is the
belief that Knowing-That (the sudden realisation of emptiness)
automatically generates Knowing-How (the dispositional capacity to live
without attachment). This is demonstrably false. A person who wishes to
be healthy may fully grasp and believe the proposition that ``smoking is
unhealthy'' (Knowing-That) while completely lacking the dispositional
capacity to actually quit (Knowing-How).

Similarly, a practitioner may realise that the ego is an illusion yet
still react to the world with ingrained habits of anger and greed
because their psychological stream---or the causal momentum of Relation
R---has not been properly conditioned. Thus, Shenxiu's gradual method is
the necessary bridge between metaphysical theory and ethical reality.
Only the gradual re-habituation of the causal stream can remove the
``dust" of conflicting desires. The Northern method of gradual
cultivation is therefore necessary to attain the dispositional nature of
enlightenment.

\section{The Pluralistic Synthesis: Narrative Logic and the Necessity
  of Insight}

If Shenxiu's method of gradual cultivation is strictly necessary to
prevent antinomianism and achieve enlightenment, one might ask: why does
the \emph{Platform Sutra} dramatise his apparent defeat? Why did the
Fifth Patriarch prioritize Huineng's verse, award him the title of the
Sixth Patriarch, and seemingly dismiss Shenxiu?

The answer lies in the epistemological sequence of enlightenment. While
Shenxiu provides the necessary method (Knowing-How), Huineng provides
the strictly necessary metaphysical target (Knowing-That). If a
practitioner engages in moral cultivation without first realising the
propositional truth of emptiness, their practice is spiritually blind
because it lacks a definitive soteriological goal. Without the orienting
framework of the reductionist-monist metaphysics---which explicitly
directs the practitioner toward absolute detachment from worldly desire---the advice to cultivate one's moral character degrades into mere
conventional morality. This is why without Huineng's ontological
assertion, Shenxiu's mandate to ``diligently polish" the mind is
incomplete: it reads as generic ethical advice on how to be a ``good"
person within the illusion, rather than a radical method for escaping
it. Because such practitioners are still attached to the concept of a
persistent self achieving a worldly goal, they remain fundamentally
trapped within \emph{samsara}.

This clarifies the narrative logic of the \emph{Platform Sutra} and
justifies the Fifth Patriarch's prioritisation of Huineng. Huineng's
verse was chosen not because it represented the entirety of
enlightenment, but because it represented the indispensable metaphysical
target. His sudden, propositional destruction of the ``Mirror Stand" acts
as a vital safeguard. It ensures that when a practitioner engages in
cultivation, they are regulating a causal stream (Relation R) rather
than feeding an ego.

It is here that we can expose the root of the traditional
interpretation's error. The orthodox commentators saw the Fifth
Patriarch rightfully crowning Huineng because his ontological insight
was strictly necessary to transcend conventional morality. However, they
made a fatal logical leap: they conflated this necessity with
sufficiency. They weaponised Huineng's narrative victory to assume that
because sudden insight was the mandatory first step, the subsequent
gradual cultivation of the causal stream was entirely obsolete.

However, the text of the \emph{Platform Sutra} itself resists this claim
of sufficiency, demonstrating that once the metaphysical target is
acquired, Shenxiu's gradualist ideas are quietly adopted. As discussed
in Section 2, Huineng actively confers the ``Formless Repentance" upon
his disciples, demanding a continuous, ongoing vigilance to reform past
and future karmic defilements. Furthermore, after receiving\\
``sudden insight,'' the Sutra depicts Huineng living in obscurity among
hunters for over a decade as a means of cultivating his dispositions.

Ultimately, the \emph{Platform Sutra} should not be read as the victory
of one mutually exclusive set of doctrines over another, but as
espousing a pluralistic theory of enlightenment. The nature of
enlightenment must be defined as both propositional insight into
emptiness and the dispositional purity to live according to it.
Consequently, the method of attainment strictly requires both sudden
awakening (to destroy the ego and secure the metaphysical target) and
gradual cultivation (to purify the causal stream).

\subsection{Historical Precedents for Methodological Pluralism}

The necessity of methodological pluralism was explicitly recognised by
the Tang Dynasty scholar-monk Guifeng Zongmi. As a patriarch of both the
Huayan school and the Heze lineage of Southern Chan, Zongmi was uniquely
positioned to diagnose the sectarian failures of his era. He observed
firsthand that the strict methodological monism of the radical Southern
School was actively leading to the exact antinomianism warned about in
this paper.

To resolve this crisis, Zongmi argued for a definitive synthesis of
Huineng and Shenxiu's views, formalising the doctrine of ``Sudden
Awakening followed by Gradual Cultivation'' (dùn wù jiàn xiū, \begin{CJK*}{UTF8}{gbsn}顿悟渐修\end{CJK*}).
Zongmi recognised that sudden insight into the ontological reality of
the self (Huineng's monism) is not the culmination of the Buddhist path,
but merely its prerequisite.

To illustrate how these two seemingly opposed methodologies must
seamlessly integrate, Zongmi utilised a series of powerful
phenomenological metaphors in his \emph{Preface to the Collected
Writings on the Source of Chan}. His most famous case study compares
enlightenment to the birth and maturation of a child:

\begin{quote}
    ``It is like a child being born all at once. At the moment of birth, its
limbs and faculties are completely formed, no different from those of an
adult. However, its physical strength is not yet mature. It requires
many months and years of gradual nourishment before it can eventually
become a fully functioning adult.''\footnote{Peter N. Gregory,
  \emph{Tsung-mi and the Sinification of Buddhism} (Princeton: Princeton
  University Press, 1991), 228.}
\end{quote}

\noindent This metaphor directly maps onto our previous analysis. The ``sudden
birth'' represents Huineng's propositional realisation of the
Metaphysical Subject: the sudden, occurrent understanding that one's
fundamental nature is entirely identical to the Buddha-nature.
Ontologically, the infant is fully human, just as the practitioner is
fully Buddha.

However, Zongmi explicitly warns that this ontological completeness does
not translate into immediate practical capacity. The ``gradual
nourishment'' represents Shenxiu's pedagogy. Even after the ``Further
Fact'' of the self is destroyed conceptually, the practitioner's
psychological continuum including their behavioural dispositions and
cognitive habits remains intact. Zongmi utilises a second metaphor to
explain this lingering causality:

\begin{quote}
    ``It is like the wind blowing across a pond. The wind may suddenly
cease, but the waves subside only gradually. The wind of ignorance may
suddenly be extinguished by awakening, but the waves of habitual
defilements can only be eliminated through gradual
cultivation.''\footnote{Gregory, \emph{Tsung-mi and the Sinification of
  Buddhism}, 228.}
\end{quote}

\noindent In Parfitian terms, Zongmi is illustrating the independence of Relation
R. The sudden cessation of the ``wind'' is the conceptual destruction of
the Cartesian Ego. But the ``waves,'' which represent the causally
linked mental states of anger, greed or selfishness, will continue to
linger. If practitioners do not try to actively calm the waves after the
cessation of the wind, they fall victim to antinomianism.

By collapsing the false narrative of mutual exclusivity, we can
therefore conclude that a Chan conception of enlightenment is neither
exclusively an occurrent realisation nor merely a dispositional state;
it is the necessary, pluralistic integration of both. It requires both
the sudden, conceptual destruction of the ``Self'' and the gradual
regulation of one's causal stream.

\section{Conclusion}

The traditional interpretation of the \emph{Platform Sutra} commits a
fundamental categorical error: it conflates the necessity of Huineng's
metaphysical insight with soteriological sufficiency. By weaponising his
narrative victory, the orthodoxy dismisses Shenxiu's teachings as an
elementary ontological mistake, thereby enforcing a dogmatic
methodological monism. As this paper has demonstrated, this orthodox
view is philosophically untenable. Regarding the nature of
enlightenment, treating propositional knowledge of emptiness as
sufficient invariably invites antinomianism. Consequently, dispositional
purity must be recognised as a necessary constituent of enlightenment.
Furthermore, through the lens of Parfitian reductionism, Shenxiu's verse
is exonerated from accusations of non-reductionism; his mandate to
``polish the mirror" is not the reification of a Cartesian Ego, but the
necessary pedagogical stewardship of the psychological continuum
(Relation R).

Regarding the method of enlightenment, Ryle's epistemology reveals that
sudden, propositional insight (Knowing-That) cannot independently
generate dispositional purity (Knowing-How). Therefore, Shenxiu's method
of gradual cultivation remains necessary to convert propositional
understanding into ethical reality. Ultimately, a tenable Chan
soteriology based on the \emph{Platform Sutra} requires the pluralistic
synthesis historically proposed by Guifeng Zongmi. The \emph{Platform
Sutra} should not be read as a contest between mutually exclusive
doctrines, but as an exposition of a plurality of doctrines required for
enlightenment. By rejecting the traditional interpretation's false
dichotomy, this analysis reveals that Chan Buddhism is not merely a
mysticism of sudden insight, but a sophisticated framework of
reductionist virtue ethics.

\refsection

\begin{hangparas}{\hangingindent}{1}
  Gregory, Peter N. \emph{Tsung-mi and the Sinification of Buddhism}.
  Princeton: Princeton University Press, 1991.


  Hume, David. \emph{A Treatise of Human Nature}. Oxford: Clarendon
  Press, 1888.


  McRae, John R. \emph{The Northern School and the Formation of Early
  Ch'an Buddhism}. Honolulu: University of Hawaii Press, 1986.


  McRae, John R. \emph{Seeing through Zen: Encounter, Transformation,
  and Genealogy in Chinese Chan Buddhism}. Berkeley: University of
  California Press, 2003.


  Parfit, Derek. \emph{Reasons and Persons}. Oxford: Oxford University
  Press, 1984.


  Ryle, Gilbert. \emph{The Concept of Mind}. London: Hutchinson, 1949.


  Siderits, Mark. \emph{Buddhism as Philosophy: An Introduction}.
  Indianapolis: Hackett Publishing, 2007.


  Yampolsky, Philip B., trans. \emph{The Platform Sutra of the Sixth
  Patriarch}. New York: Columbia University Press, 1967.


  Zongmi, Guifeng. \emph{Chanyuan zhuquanji duxu} {[}Preface to the
  Collected Writings on the Source of Chan{]}. Translated and quoted in
  Peter N. Gregory, \emph{Tsung-mi and the Sinification of Buddhism}.
  Princeton: Princeton University Press, 1991.
\end{hangparas}